\documentclass[11pt,a4paper]{article}

% ---------------------------------------------------------
% PREAMBLE (stable, warning-free, Overleaf-safe)
% ---------------------------------------------------------
\usepackage[utf8]{inputenc}
\usepackage[T1]{fontenc}
\usepackage{lmodern}
\usepackage{amsmath, amssymb, amsfonts}
\usepackage{bm}
\usepackage{geometry}
\usepackage{graphicx}
\usepackage{hyperref}
\usepackage{cite}
\usepackage{authblk}
\usepackage{physics}
\usepackage{mathtools}
\usepackage{textgreek}
\geometry{margin=2.4cm}

\title{\textbf{Companion LX: Weak-Lensing Peak Statistics in the Relaxation-Driven Cyclic Cosmology}}
\author{Michael Lehmann \\ \texttt{mi.lehmann@gmx.de}}
\date{April 2026}

\begin{document}
\maketitle

\begin{abstract}
Weak-lensing peak statistics probe the nonlinear structure of the gravitational 
field by identifying local maxima in convergence maps. In the standard 
cosmological model, the abundance and morphology of peaks are determined by the 
nonlinear matter distribution, the halo mass function, and the projection of 
large-scale structure. In the Relaxation-Driven Cyclic Cosmology (RDCC), the 
scale-dependent evolution of the gravitational potential modifies the depth, 
coherence, and spatial distribution of peaks in a characteristic way. This 
Companion develops a continuous description of weak-lensing peak statistics in 
RDCC, deriving how the relaxation scale influences the peak abundance, the 
high-κ tail, and the morphology of nonlinear structures. The resulting framework 
provides a distinctive probe of the RDCC gravitational field beyond two- and 
three-point statistics.
\end{abstract}
% ---------------------------------------------------------
\section{Introduction}

Weak-lensing peak statistics identify local maxima in convergence maps and provide 
a powerful probe of nonlinear structure formation. High peaks correspond to massive 
halos or line-of-sight projections, while low peaks reflect the interplay of 
nonlinear structure and large-scale modes.

In the standard cosmological model, the peak abundance and morphology are 
determined by the nonlinear matter distribution, the halo mass function, and the 
projection of large-scale structure. In RDCC, the gravitational potential evolves 
differently. The relaxation mechanism suppresses the decay of small-scale modes and 
enhances the coherence of large-scale modes. This modifies the depth, coherence, 
and spatial distribution of peaks in a scale-dependent manner, producing a 
distinctive signature in weak-lensing peak statistics.
% ---------------------------------------------------------
\section{The RDCC Convergence Field and Peak Definition}

The convergence field is
\begin{equation}
    \kappa(\hat{\mathbf{n}})
    = \int_{0}^{\chi_{H}}
      d\chi\,
      W(\chi)\,
      \delta_{m}\left(\chi,\hat{\mathbf{n}}\right).
\end{equation}

In RDCC, the matter density evolves as
\begin{equation}
    \delta_{m,\mathrm{RDCC}}(k,z)
    = \delta_{m}(k,z)
      \left[
        1 - \chi_{\delta}
        \left(\frac{k}{k_{\mathrm{rel}}}\right)^{\alpha}
      \right].
\end{equation}

The RDCC convergence becomes
\begin{equation}
    \kappa_{\mathrm{RDCC}}(\ell)
    = \kappa(\ell)
      \left[
        1 - \chi_{\kappa}
        \left(\frac{\ell}{\ell_{\mathrm{rel}}}\right)^{\alpha}
      \right].
\end{equation}

A peak is defined as a pixel where
\begin{equation}
    \kappa(\hat{\mathbf{n}}) > \kappa(\hat{\mathbf{n}} + \Delta\hat{\mathbf{n}}),
\end{equation}
for all neighboring directions $\Delta\hat{\mathbf{n}}$.
% ---------------------------------------------------------
\section{Peak Abundance in RDCC}

The peak abundance is given by
\begin{equation}
    n_{\mathrm{peak}}(\nu)
    = \frac{dN_{\mathrm{peak}}}{d\nu},
\end{equation}
where $\nu = \kappa / \sigma_{\kappa}$.

In RDCC, the variance becomes
\begin{equation}
    \sigma_{\kappa,\mathrm{RDCC}}^{2}
    = \sigma_{\kappa}^{2}
      \left[
        1 - \chi_{\sigma}
        \left(\frac{\ell}{\ell_{\mathrm{rel}}}\right)^{\alpha}
      \right].
\end{equation}

The peak abundance becomes
\begin{equation}
    n_{\mathrm{peak,RDCC}}(\nu)
    = n_{\mathrm{peak}}(\nu)
      \left[
        1 - \chi_{n}
        \left(\frac{\ell}{\ell_{\mathrm{rel}}}\right)^{\alpha}
      \right].
\end{equation}

This modification suppresses small-scale peaks and enhances large-scale peaks.
% ---------------------------------------------------------
\section{High-Peak Tail and RDCC Signatures}

High peaks correspond to massive halos or strong line-of-sight projections. The 
high-κ tail is given by
\begin{equation}
    n_{\mathrm{peak}}(\nu \gg 1)
    \propto \exp\left(-\frac{\nu^{2}}{2}\right).
\end{equation}

In RDCC, the high-κ tail becomes
\begin{equation}
    n_{\mathrm{peak,RDCC}}(\nu \gg 1)
    \propto
    \exp\left[
      -\frac{\nu^{2}}{2}
      \left(
        1 - \chi_{\mathrm{tail}}
        \left(\frac{\ell}{\ell_{\mathrm{rel}}}\right)^{\alpha}
      \right)
    \right].
\end{equation}

This modification reflects the relaxation-driven suppression of small-scale 
potential wells.
% ---------------------------------------------------------
\section{Peak Morphology and RDCC Signatures}

The relaxation mechanism enhances the coherence of large-scale modes. Peak 
morphology therefore exhibits:

\begin{itemize}
    \item broader, smoother peaks at large scales,
    \item suppressed small-scale substructure,
    \item a characteristic turnover in the peak-size distribution near $\ell_{\mathrm{rel}}$,
    \item modified correlations between peak height and curvature.
\end{itemize}

These signatures provide a direct observational probe of the RDCC gravitational 
field in real-space mass maps.
% ---------------------------------------------------------
\section{Conclusion}

Weak-lensing peak statistics in the Relaxation-Driven Cyclic Cosmology reflect the 
scale-dependent evolution of the gravitational potential and the nonlinear matter 
distribution. The relaxation mechanism modifies the peak abundance, the high-κ 
tail, and the morphology of peaks in a scale-dependent manner, producing 
distinctive signatures in weak-lensing surveys. These effects provide a direct 
observational window into the relaxation scale and the temporal structure of the 
RDCC gravitational field. The present Companion establishes the theoretical 
foundation for future studies of nonlinear structure, mass mapping, and the 
interplay between light and gravity in RDCC.

\bibliographystyle{apsrev4-2}
\nocite{*}
\bibliography{references}

\end{document}
